\begin{ex}
Một khối lập phương có bao nhiêu mặt?
\shortans{6}
\loigiai{
Khối lập phương có đúng 6 mặt phẳng bằng nhau và đều là hình vuông.
\textbf{Đáp số:} $6$.
}
\end{ex}

\begin{ex}
Một khối hộp chữ nhật có bao nhiêu mặt?
\shortans{6}
\loigiai{
Khối hộp chữ nhật có đúng 6 mặt (các mặt đều là hình chữ nhật hoặc hình vuông).
\textbf{Đáp số:} $6$.
}
\end{ex}

\begin{ex}
Một khối lập phương có bao nhiêu đỉnh (góc nhọn)?
\shortans{8}
\loigiai{
Khối lập phương có 8 đỉnh (4 đỉnh mặt trên và 4 đỉnh mặt dưới).
\textbf{Đáp số:} $8$.
}
\end{ex}

\begin{ex}
Một khối hộp chữ nhật có bao nhiêu đỉnh?
\shortans{8}
\loigiai{
Khối hộp chữ nhật có 8 đỉnh.
\textbf{Đáp số:} $8$.
}
\end{ex}

\begin{ex}
Các mặt của một khối lập phương đều là hình gì? (1: Hình tròn; 2: Hình tam giác; 3: Hình vuông; 4: Hình chữ nhật).
\shortans{3}
\loigiai{
Tất cả 6 mặt của khối lập phương đều là hình vuông (ghi số 3).
\textbf{Đáp số:} $3$.
}
\end{ex}

\begin{ex}
Hai khối lập phương có tất cả bao nhiêu mặt?
\shortans{12}
\loigiai{
Mỗi khối lập phương có 6 mặt. Hai khối có: $6 + 6 = 12$ (mặt).
\textbf{Đáp số:} $12$.
}
\end{ex}

\begin{ex}
Một khối lập phương đặt trên mặt bàn. Hỏi có bao nhiêu mặt của khối lập phương không tiếp xúc với mặt bàn?
\shortans{5}
\loigiai{
Có 1 mặt đáy tiếp xúc với mặt bàn. Số mặt không tiếp xúc là: $6 - 1 = 5$ (mặt).
\textbf{Đáp số:} $5$.
}
\end{ex}

\begin{ex}
Một khối hộp chữ nhật đặt trên sàn nhà. Có bao nhiêu mặt đáy của khối hộp chữ nhật tiếp xúc với sàn nhà?
\shortans{1}
\loigiai{
Chỉ có đúng $1$ mặt đáy tiếp xúc với sàn nhà.
\textbf{Đáp số:} $1$.
}
\end{ex}

\begin{ex}
Bạn Nam có 3 khối lập phương và 4 khối hộp chữ nhật. Hỏi Nam có tất cả bao nhiêu khối hình học?
\shortans{7}
\loigiai{
Tổng số khối hình Nam có là: $3 + 4 = 7$ (khối hình).
\textbf{Đáp số:} $7$.
}
\end{ex}

\begin{ex}
Trong một hộp đồ chơi có 10 khối hình gồm khối lập phương và khối hộp chữ nhật. Biết có 6 khối lập phương, hỏi có bao nhiêu khối hộp chữ nhật?
\shortans{4}
\loigiai{
Số khối hộp chữ nhật là: $10 - 6 = 4$ (khối).
\textbf{Đáp số:} $4$.
}
\end{ex}

\begin{ex}
Một khối lập phương có bao nhiêu cạnh?
\shortans{12}
\loigiai{
Khối lập phương có 12 cạnh (4 cạnh trên, 4 cạnh dưới và 4 cạnh bên).
\textbf{Đáp số:} $12$.
}
\end{ex}

\begin{ex}
Một khối hộp chữ nhật có bao nhiêu cạnh?
\shortans{12}
\loigiai{
Khối hộp chữ nhật có đúng 12 cạnh.
\textbf{Đáp số:} $12$.
}
\end{ex}

\begin{ex}
Hiệu giữa số mặt của một khối lập phương và một khối hộp chữ nhật là bao nhiêu?
\shortans{0}
\loigiai{
Cả hai khối đều có 6 mặt, hiệu là: $6 - 6 = 0$.
\textbf{Đáp số:} $0$.
}
\end{ex}

\begin{ex}
Bác thợ mộc có 8 khối gỗ hình hộp chữ nhật. Bác đã dùng 5 khối gỗ để đóng bàn ghế. Hỏi bác thợ còn lại bao nhiêu khối gỗ?
\shortans{3}
\loigiai{
Số khối gỗ còn lại là: $8 - 5 = 3$ (khối gỗ).
\textbf{Đáp số:} $3$.
}
\end{ex}

\begin{ex}
Xếp 2 khối lập phương nhỏ cạnh nhau tạo thành một khối hình mới. Khối hình mới đó là hình gì? (1: Khối cầu; 2: Khối trụ; 3: Khối hộp chữ nhật; 4: Khối lập phương).
\shortans{3}
\loigiai{
Ghép 2 khối lập phương cạnh nhau sẽ tạo thành một khối hộp chữ nhật (ghi số 3).
\textbf{Đáp số:} $3$.
}
\end{ex}

\begin{ex}
Cần bao nhiêu khối lập phương nhỏ bằng nhau để xếp thành một khối lập phương lớn hơn gồm 2 tầng, mỗi tầng có 4 khối?
\shortans{8}
\loigiai{
Số khối lập phương nhỏ cần dùng là: $4 + 4 = 8$ (khối).
\textbf{Đáp số:} $8$.
}
\end{ex}

\begin{ex}
Bé dùng 5 khối lập phương xếp chồng lên nhau thành một cột đứng. Bé nhấc bớt đi 2 khối ở trên cùng. Hỏi cột đó còn lại bao nhiêu khối lập phương?
\shortans{3}
\loigiai{
Số khối lập phương còn lại là: $5 - 2 = 3$ (khối).
\textbf{Đáp số:} $3$.
}
\end{ex}

\begin{ex}
Bao diêm thường có dạng khối hình gì? (1: Khối lập phương; 2: Khối hộp chữ nhật; 3: Khối trụ; 4: Khối cầu).
\shortans{2}
\loigiai{
Bao diêm có dạng khối hộp chữ nhật (ghi số 2).
\textbf{Đáp số:} $2$.
}
\end{ex}

\begin{ex}
Viên xúc xắc đổ súc sắc thường có dạng khối hình gì? (1: Khối lập phương; 2: Khối hộp chữ nhật; 3: Khối trụ; 4: Khối cầu).
\shortans{1}
\loigiai{
Viên xúc xắc có 6 mặt vuông bằng nhau nên có dạng khối lập phương (ghi số 1).
\textbf{Đáp số:} $1$.
}
\end{ex}

\begin{ex}
Tủ lạnh gia đình thường có dạng hình khối gì? (1: Khối lập phương; 2: Khối hộp chữ nhật; 3: Khối trụ; 4: Khối cầu).
\shortans{2}
\loigiai{
Tủ lạnh gia đình có dạng khối hộp chữ nhật đứng (ghi số 2).
\textbf{Đáp số:} $2$.
}
\end{ex}

\begin{ex}
Trên bàn có 2 khối lập phương và số khối hộp chữ nhật nhiều hơn khối lập phương là 3 khối. Hỏi trên bàn có bao nhiêu khối hộp chữ nhật?
\shortans{5}
\loigiai{
Số khối hộp chữ nhật là: $2 + 3 = 5$ (khối).
\textbf{Đáp số:} $5$.
}
\end{ex}

\begin{ex}
Tổng số đỉnh của một khối lập phương và một khối hộp chữ nhật là bao nhiêu đỉnh?
\shortans{16}
\loigiai{
Mỗi khối có 8 đỉnh. Tổng số đỉnh của cả hai khối là: $8 + 8 = 16$ (đỉnh).
\textbf{Đáp số:} $16$.
}
\end{ex}

\begin{ex}
Xếp 4 khối lập phương thành một hàng ngang. Khối hình đó có tất cả bao nhiêu khối lập phương?
\shortans{4}
\loigiai{
Khối hình đó được tạo từ đúng $4$ khối lập phương nhỏ.
\textbf{Đáp số:} $4$.
}
\end{ex}

\begin{ex}
Có 10 khối gỗ hình lập phương, em dùng 4 khối xếp tòa tháp A, dùng 3 khối xếp tòa tháp B. Hỏi em còn lại bao nhiêu khối gỗ?
\shortans{3}
\loigiai{
Tổng số khối gỗ đã dùng là: $4 + 3 = 7$ (khối).\nSố khối gỗ còn lại là: $10 - 7 = 3$ (khối gỗ).
\textbf{Đáp số:} $3$.
}
\end{ex}

\begin{ex}
Một chiếc hộp sữa bột có 6 mặt chữ nhật. Chiếc hộp đó có dạng khối gì? (Khối hộp chữ nhật: ghi 2; Khối lập phương: ghi 1).
\shortans{2}
\loigiai{
Chiếc hộp đó có dạng khối hộp chữ nhật (ghi số 2).
\textbf{Đáp số:} $2$.
}
\end{ex}
